% Section 2.4, from lectures/L02/appendix/c_exercises.md.

\section{Övningar}
\label{c2:sec:exercises}\label{c2:app:c}

\begin{aside}
\textbf{Så kontrollerar du ditt arbete.} Varje uttryck du förenklar kan kontrolleras med en
sanningstabell: räkna ut det ursprungliga och det förenklade uttrycket för alla kombinationer av
ingångarna, och jämför kolumnerna. Är de lika på varje rad är förenklingen rätt, oavsett vilken väg
du tog dit. Varje grindnät kan dessutom byggas i CircuitVerse och kontrolleras rad för rad, enligt
\secref{c2:a9}.

Lösningsförslagen finns i bilaga~\ref{app:answers}.

Varje övning är märkt med sin sort. Övning~\ref{c2:ex:7} är kapitlets
\textbf{kontrolluppgift}.
\end{aside}

\exercise{Grindar med tre ingångar}{Förståelse}
\label{c2:ex:1}
\begin{parts}
  \item Skriv sanningstabellen för en AND-grind, en OR-grind och en XOR-grind med tre ingångar A,
    B och C, i samma tabell.
  \item Beskriv med en mening var för varje grind när utgången är \code{1}.
  \item På hur många av de åtta raderna är AND-grinden \code{1}? OR-grinden? Förklara varför
    svaren är så olika.
\end{parts}

\exercise{Från ord till uttryck}{Förståelse}
\label{c2:ex:2}
Skriv ett booleskt uttryck för varje beskrivning. Välj själv bokstäver för signalerna, och ange vad
de betyder.
\begin{parts}
  \item Larmet ska ljuda om dörren är öppen och larmet är aktiverat.
  \item Pumpen ska gå om nivån i tanken är låg och ingen har tryckt på det manuella stoppet.
  \item Varningslampan ska lysa om minst en av tre dörrar är öppen.
  \item Beskriv med ord vad \mbox{\code{X = AB' + C}} betyder, om A är
    \enquote{maskinen är påslagen}, B är \enquote{skyddet är på plats} och C är
    \enquote{nödlarmet är utlöst}.
  \item Beräkna \mbox{\code{AB + C}} och \mbox{\code{A(B + C)}} för A = 0, B = 0 och C = 1. Vad
    visar resultatet om prioritetsordningen?
\end{parts}

\exercise{Räknelagarna}{Räkna för hand}
\label{c2:ex:3}
Förenkla med lagarna i \secref{c2:a4}, och ange vilken lag du använder.
\begin{parts}
  \item \code{A + A}
  \item \code{A · A'}
  \item \code{A + 1}
  \item \code{(A')'}
  \item \code{A + AB}
  \item \code{A(A + B)}
  \item \code{A + A'B}
  \item \code{AB + A'B}
\end{parts}

\exercise{De Morgan}{Räkna för hand}
\label{c2:ex:4}
\begin{parts}
  \item Skriv om \code{(AB)'} utan någon inversion av ett helt uttryck.
  \item Skriv om \mbox{\code{(A + B + C)'}} på samma sätt.
  \item Förenkla \code{(A'B)'}. Tänk på att \mbox{\code{(A')' = A}}.
  \item Förenkla \mbox{\code{(A + B')'}}.
  \item Visa med en sanningstabell att \mbox{\code{(AB)' = A' + B'}}.
\end{parts}

\exercise{Från sanningstabell till uttryck}{Räkna för hand}
\label{c2:ex:5}

\begin{narrowtable}{@{}cccc@{}}
  \thead{A} & \thead{B} & \thead{C} & \thead{X} \\
  \midrule
  0 & 0 & 0 & 1 \\
  0 & 0 & 1 & 0 \\
  0 & 1 & 0 & 1 \\
  0 & 1 & 1 & 0 \\
  1 & 0 & 0 & 0 \\
  1 & 0 & 1 & 0 \\
  1 & 1 & 0 & 1 \\
  1 & 1 & 1 & 1 \\
\end{narrowtable}

\begin{parts}
  \item Läs av uttrycket för X på SP-form, med en minterm per etta.
  \item Förenkla uttrycket algebraiskt. Du ska få två termer med två variabler var.
  \item Kontrollera ditt förenklade uttryck mot tabellen, rad för rad.
  \item Hur många grindar behöver det oförenklade respektive det förenklade uttrycket, om du har
    NOT-grindar och AND- och OR-grindar med så många ingångar du vill?
\end{parts}

\exercise{Algebraisk förenkling}{Räkna för hand}
\label{c2:ex:6}
Förenkla så långt det går, och kontrollera varje svar med en sanningstabell.
\begin{parts}
  \item \code{AB + AB' + A'B}
  \item \code{(A + B)(A + B')}
  \item \mbox{\code{S'R + SR' + SR}}, som behövdes i härledningen av bromsassistenten i
    \secref{c2:a10}. Visa att det är \mbox{\code{S + R}}.
  \item \code{A'B'C + A'BC + ABC}
\end{parts}

\exercise{Kontroll: analys av ett grindnät}{Kontroll}
\label{c2:ex:7}
\emph{Analysera för hand, bygg och simulera, jämför.}

\bookfigure{l02_ex_network}{Grindnätet i övning~\ref{c2:ex:7}.}{c2:fig:ex_network}

\noindent Gör delarna i ordning, och skriv ned varje svar innan du går vidare.

\task{a) För hand}
Ge NOR-grindens utgång namnet P och AND-grindens utgång namnet Q. Skriv uttrycken för P, Q och X.

\task{b)}
Räkna fram sanningstabellen med kolumnerna A, B, C, P, Q och X, som i \secref{c2:a8}.

\task{c) I CircuitVerse}
Bygg nätet precis som det är ritat, och gå igenom alla åtta kombinationerna. Stämmer simuleringen
med din tabell?

\task{d)}
Om inte: är det tabellen eller bygget som är fel? Kontrollera den rad som skiljer genom att räkna ut
P och Q för just den raden en gång till, och genom att följa varje ledning i bygget.

\task{e)}
Beskriv med en mening när X är \code{1}.

\exercise{Från uttryck till grindnät}{Konstruktion}
\label{c2:ex:8}
\begin{parts}
  \item Rita grindnätet för \mbox{\code{X = A'B + BC'}} med NOT-, AND- och OR-grindar. Hur många
    grindar går det åt?
  \item Bryt ut B ur uttrycket, och använd De Morgan på det som blir kvar. Visa att samma funktion
    går att bygga med en NAND-grind och en AND-grind.
  \item Rita grindnätet för \mbox{\code{Y = (A + B)C'}}. Hur många grindar går det åt?
  \item Bygg nätet från b) i CircuitVerse och kontrollera det mot en sanningstabell för
    \mbox{\code{A'B + BC'}}.
\end{parts}

\exercise{Allt av NAND}{Konstruktion}
\label{c2:ex:9}
\Secref{c2:a5} påstår att varje funktion går att bygga med enbart NAND-grindar. Visa det för de tre
grundoperationerna.
\begin{parts}
  \item Visa algebraiskt att en NAND-grind med båda ingångarna ihopkopplade är en NOT-grind.
  \item Hur bygger du en AND av NAND-grindar? Hur många behövs?
  \item Utgå från \mbox{\code{A + B = ((A + B)')'}} och använd De Morgan för att skriva OR med
    enbart NAND-operationer. Hur många NAND-grindar behövs?
  \item Bygg alla tre i CircuitVerse med enbart NAND-grindar, och kontrollera var och en mot sin
    sanningstabell.
\end{parts}

\exercise{Från uttryck till kontaktnät}{Konstruktion}
\label{c2:ex:10}
Rita ett strömvägsschema mellan +24~V och 0~V för varje uttryck, med tryckknappar och lampan H1.
Ange för varje kontakt om den är slutande eller brytande.
\begin{parts}
  \item \code{H1 = S1 · S2'}
  \item \code{H1 = S1 + S2 · S3}
  \item \code{H1 = (S1 + S2) · S3'}
  \item Hur många kontaktelement behöver varje knapp ha i a)--c)? Och hur många i XOR-kopplingen i
    \secref{c2:b5}?
\end{parts}

\exercise{Analys av kontaktnät}{Räkna för hand}
\label{c2:ex:11}

\bookfigure{l02_ex_contact}{Kontaktnäten a) och b) i övning~\ref{c2:ex:11}.}{c2:fig:ex_contact}

\begin{parts}
  \item Skriv uttrycket för H1 i nät a), och ta fram sanningstabellen för S1, S2 och S3.
  \item Gör samma sak för H2 i nät b).
  \item I nät b): vilka kombinationer tänder H2 utan att S1 eller S2 är nedtryckt? Förklara med
    ritningen varför.
\end{parts}

\exercise{Självhållning}{Förståelse}
\label{c2:ex:12}
Utgå från självhållningskretsen i \secref{c2:b7}, figur~\ref{c2:fig:relay_self_hold}.
\begin{parts}
  \item Beskriv vad som händer med spolen K1 och lampan H1 när du trycker på S1, och sedan när du
    släpper den.
  \item Beskriv vad som händer när du sedan trycker på S2, och när du släpper den.
  \item Vad händer om någon trycker på S1 och S2 samtidigt? Vilken av knapparna \enquote{vinner}?
  \item Vid monteringen kopplas stoppknappen S2 av misstag med sitt slutande kontaktelement i
    stället för sitt brytande. Beskriv hur maskinen beter sig.
  \item Förklara varför kretsen inte är ett kombinatoriskt nät. Använd ordet minne.
\end{parts}

\exercise{Bromsassistenten, felkopplad}{Räkna för hand, Fördjupning}
\label{c2:ex:13}
En konstruktör vill spara en grind och kopplar bromsassistenten från \secref{c2:a10} som

\begin{console}
B = (D + S + R)F'
\end{console}

\noindent Det är en OR-grind med tre ingångar följd av en AND med den inverterade felsignalen, och
det ser rimligt ut.
\begin{parts}
  \item Ta fram sanningstabellen för det nya uttrycket, och jämför den rad för rad med tabellen i
    \secref{c2:a10}. På vilka rader skiljer de sig?
  \item Vad har de raderna gemensamt? Beskriv vad som händer fysiskt i bilen i en sådan situation.
  \item Vad säger det här om värdet av en sanningstabell som är härledd ur kravet, jämfört med en
    som är avläst ur en krets man redan har byggt?
\end{parts}
